% Historical endpoint changes used only to illustrate the index scale.
% Values are computed from the complete 193-country, 1960--2024 panel.
\begin{table}[!htb]
\centering
\caption{Historical Realignments on the Index Scale}
\label{tab:episodes_1sd}
\footnotesize
\setlength{\tabcolsep}{4pt}
\resizebox{\linewidth}{!}{%
\begin{tabular}{lclc}
\toprule
Benchmark or episode & $\Delta$ index (SD) & Benchmark or episode & $\Delta$ index (SD) \\
\midrule
Interquartile shift (25th$\to$75th) & 0.186 (1.32) & Iraq: Gulf War and sanctions (1989--94) & $-0.58$ ($-4.1$) \\
South Africa: apartheid ends (1986--96) & $+0.69$ ($+4.9$) & Vietnam: normalization with the US (1990--2000) & $+0.19$ ($+1.4$) \\
Iran: Islamic Revolution (1977--83) & $-0.48$ ($-3.4$) & Russia: deterioration through Crimea (2005--15) & $-0.14$ ($-1.0$) \\
Poland: post-communist transition (1988--98) & $+0.27$ ($+1.9$) & Indonesia: post-Suharto reorientation (1996--2004) & $+0.13$ ($+1.0$) \\
China: reform and opening (1976--86) & $+0.15$ ($+1.1$) & Egypt: Camp David-era realignment (1973--81) & $+0.05$ ($+0.3$) \\
\bottomrule
\end{tabular}}
\note{$\Delta$ is the end-year minus start-year value for each selected window. Parentheses report the change in standard deviations, using the full-sample pooled $\sigma=0.141$. The interquartile benchmark is the shift from the 25th percentile (0.154) to the 75th percentile (0.340). Episode labels provide historical context, and examples are chosen for recognizability rather than as a ranking.}
\end{table}
